% !PRINTED 173
\weberpage{173}
%<<CONT
$\varkappa_\nu$ approaches the limit zero as $\nu$ increases infinitely. From this one obtains for $q$ the representation:

\weq{22}{q = \mathop{\mathrm{Lim}}\limits_{\nu=\infty} \sqrt[2^{\nu-1}]{\frac{\varkappa_\nu}{4\varkappa_\nu'}} = \mathrm{Lim} \sqrt[2^{\nu-1}]{\frac{\varkappa_\nu}{4}},}

which is suited to computation for real values of $\varkappa$.

From each of the formulae (21) one can derive a series development of $q$ in ascending powers of $\varkappa_\nu$, whose convergence increases as $\nu$ increases. If in particular one applies formula (20), one obtains a development which was communicated by \gesperrt{Weierstrass} in the Monatsberichte der Berliner Akademie for the year 1883, and whose domain of convergence fills the plane $\varkappa^2$ twice.

The coefficients of this development are most simply computed from:

\weq{23}{\sqrt{\varkappa_2} = \frac{\vartheta_{10}(0,4\omega)}{\vartheta_{00}(0,4\omega)} = \frac{2q+2q^9+2q^{25}+\cdots}{1+2q^4+2q^{16}+\cdots},}

and one thus obtains, by the method of undetermined coefficients:

\weq{24}{q = \frac{\sqrt{\varkappa_2}}{2} + 2\left(\frac{\sqrt{\varkappa_2}}{2}\right)^5 + 15\left(\frac{\sqrt{\varkappa_2}}{2}\right)^9 + 150\left(\frac{\sqrt{\varkappa_2}}{2}\right)^{13} + 1707\left(\frac{\sqrt{\varkappa_2}}{2}\right)^{17} + \cdots}

In the same way, conversely, $\sqrt{\varkappa_2}$ can also be developed by (23) into a series proceeding in powers of $q$:

\weq{25}{\tfrac{1}{2}\sqrt{\varkappa_2} = q - 2q^5 + 5q^9 - 10q^{13} + 18q^{17}\cdots,}

which one also obtains by inversion of the series (24).

With this the question posed at the beginning of this section is completely answered.

As regards the application of the formulae developed here to the numerical computation of $q$, it may still be remarked that, when $\varkappa^2$ lies nearer to 1, one obtains a better converging procedure if one interchanges $\varkappa^2$ with $\varkappa'^2$. The calculation then yields at first not $q$, but

\[
q' = e^{-\frac{\pi i}{\omega}},
\]

from which, however, one obtains $q$ from the formula:

\[
\log q \log q' = \pi^2.
\]
